\documentclass[12pt]{article}
\usepackage[utf8]{inputenc}
\usepackage[T1]{fontenc}
\usepackage[margin=1in]{geometry}
\usepackage{setspace}
\usepackage{xcolor}
\usepackage[colorlinks=true, urlcolor=blue!60!black, linkcolor=blue!60!black]{hyperref}
\setlength{\parindent}{2em}
\setlength{\parskip}{1.25em}
\renewcommand{\baselinestretch}{1.0}
\begin{document}
\begin{center}
{\LARGE \textit{Hernandez v.\ Texas}, 347 U.S.\ 475 (1954)\par}
\vspace{0.8em}
{\large Case Brief --- Four-Page Edition\par}
\end{center}
\vspace{1em}

\section{Memory Jogger}

The Equal Protection Clause is not limited to a Black--white ``two-class theory,'' and bars the systematic exclusion of Mexican Americans---or any other locally subordinated distinct class---from jury service (\href{https://supreme.justia.com/cases/federal/us/347/475/}{Justia}).

\section{Detailed Case Facts}

In 1951 Pete Hernandez, a Mexican-American agricultural worker, was indicted for murder by an all-Anglo grand jury in Jackson County, Texas, convicted by an all-Anglo petit jury, and sentenced to life. His lawyers, including Gus Garc\'{i}a and Carlos Cadena, brought the first Supreme Court case argued by Mexican-American attorneys. About 14\% of the county's residents and 11\% of its adult males had Spanish surnames, and 6\% to 7\% of freeholders were of Mexican descent. Texas stipulated that no person with a Mexican or Latin American name had served on a jury commission, grand jury, or petit jury in 25 years---out of more than 6,000 jurors---although qualified Mexican Americans lived there. Hernandez also proved local segregation: residents distinguished ``white'' from ``Mexican,'' Mexican-American children attended a segregated school for the first four grades, a restaurant posted ``No Mexicans Served,'' and the courthouse had a men's restroom marked ``Colored Men'' and ``Hombres Aqui.'' Five jury commissioners testified only that they had not discriminated.

\section{Procedural History}

The trial court denied motions to quash the indictment and the jury panel. The Texas Court of Criminal Appeals affirmed, holding that the Fourteenth Amendment recognized only two classes, white and Negro, and that Mexican Americans were white. The Supreme Court reversed on May 3, 1954, two weeks before \textit{Brown v. Board of Education}. (On retrial, Hernandez was convicted again.)

\section{Judicial Votes}

Reversed, \textbf{9--0}.
\begin{itemize}
\item \textbf{Opinion of the Court:} Chief Justice Warren, for a unanimous Court.
\item \textbf{Concurrences and dissents:} None.
\end{itemize}

\section{Holding}

The Equal Protection Clause is not confined to a two-class theory; it forbids excluding otherwise eligible persons from juries because of ancestry or national origin. A challenger must (1) prove as a fact that the group is a distinct class in the local community and (2) make a prima facie case of systematic exclusion under the \textit{Norris v. Alabama} rule of exclusion. General denials of discriminatory intent cannot rebut that case. Footnote 9 reserved whether courts may take judicial notice that Mexican Americans are a separate class, and the Court disclaimed any right to proportional representation.

\section{Analysis of Opinions}

\textbf{Majority (Warren).} Drawing on \textit{Strauder}, \textit{Truax v. Raich}, \textit{Takahashi}, and \textit{Hirabayashi}, Warren reasoned that groups needing protection change as community prejudices change; Texas's own courts had barred excluding Roman Catholics from juries. On the first step, the attitude of the community---testimony, school segregation, the restaurant sign, and the restroom---proved that Mexican Americans were a distinct class. Spanish surnames were a fair proxy because commissioners worked from lists of names, which gave ``ready identification'' as skin color did for race. On the second step, with qualified Mexican Americans in the county and none chosen in 25 years among over 6,000 jurors, mere chance could not explain the result, and the commissioners' general assertions of good faith could not overcome the prima facie proof. No Justice wrote separately.

\section{Examples: Future Applications}

\begin{enumerate}
\item \textbf{Same-Side (Violation): Hmong Americans.} In a county where Hmong Americans are 12\% of adults and face local segregation, none has been called among 4,500 jurors in 20 years. Community attitudes establish a class, and total exclusion triggers the \textit{Norris} inference.
\item \textbf{Same-Side (Violation): Key-Man System.} A facially neutral key-man system selects no Arab-American grand jurors in 15 years, and commissioners say only that they chose respectable acquaintances. A neutral statute is invalid as applied.
\item \textbf{Opposite-Side (No Violation): Random Draw.} Mexican Americans serve proportionally in the county, but by chance none sits on the defendant's jury. \textit{Hernandez} guarantees freedom from systematic exclusion, not representation on a particular jury.
\item \textbf{Opposite-Side (No Violation): Neutral Qualifications.} Uniform exclusion of felons and non-English speakers, with no group singled out, rests on reasonable classifications.
\item \textbf{Fence-Sitter: Severe Underrepresentation.} Mexican Americans are 30\% of residents but 4\% of grand jurors over a decade. Without an inexorable zero, the case is closer; \textit{Castaneda v. Partida} (1977) later accepted statistical-disparity proof.
\end{enumerate}

\section{Critique}

\textbf{Originalist and formalist views.} \textit{Slaughter-House} once doubted that equal protection reached beyond the freedmen, but textualists endorse rejecting the two-class theory because the Amendment protects ``any person.'' Formalist and colorblindness critics question Warren's open-ended sociological test and his willingness to find discrimination without a conscious decision, in tension with \textit{Washington v. Davis}.

\textbf{Progressive and LatCrit views.} Ian Haney L\'{o}pez praised \textit{Hernandez} as ``the single most insightful Supreme Court opinion on race ever handed down,'' because it treated group identity as a social construct. Yet scholars fault the class-apart strategy: by requiring proof of local community attitudes county by county, the Court imposed a repetitive evidentiary burden on Latino litigants and left the Black--white binary intact.

\textbf{Competing values.} The case pits anti-subordination against formal anti-classification, federal supervision of state criminal justice against state discretion over juries, and structural proof of exclusion against proof of individual intent.

\section{Key Quotations}

\begin{quote}
``Throughout our history, differences in race and color have defined easily identifiable groups which have at times required the aid of the courts in securing equal treatment under the laws. But community prejudices are not static, and, from time to time, other differences from the community norm may define other groups which need the same protection.''
\end{quote}
\noindent --- Chief Justice Warren, for the Court, 347 U.S. at 478.

\begin{quote}
``When the existence of a distinct class is demonstrated, and it is further shown that the laws, as written or as applied, single out that class for different treatment not based on some reasonable classification, the guarantees of the Constitution have been violated. The Fourteenth Amendment is not directed solely against discrimination due to a `two-class theory'---that is, based upon differences between `white' and Negro.''
\end{quote}
\noindent --- Chief Justice Warren, for the Court, 347 U.S. at 478.

\begin{quote}
``But it taxes our credulity to say that mere chance resulted in there being no members of this class among the over six thousand jurors called in the past 25 years. The result bespeaks discrimination, whether or not it was a conscious decision on the part of any individual jury commissioner.''
\end{quote}
\noindent --- Chief Justice Warren, for the Court, 347 U.S. at 482.

\end{document}
